\chapter{Enrichments and Quantitative Semantics}\label{ch:enrich}

Mol describes what can be composed and which behaviors are equivalent; it says
nothing about cost. This chapter enriches the structure with numbers. The
central idea is that Mol is not just a category but a \emph{resource-sensitive
monoidal category}: every molecule carries a cost vector --- execution time,
memory footprint, and cache misses --- and the composition operations of
\autoref{ch:mol} act on costs in a controlled way. Sequential composition
adds costs, tensor composition is subadditive (parallel work shares fixed
costs), and a \emph{refinement} preorder orders molecules by quality inside a
behavior class.

The chapter is the quantitative backbone of the thesis. NT25 makes the cost
enrichment well-defined: additive composition cost, the triangle inequality,
and tensor subadditivity. NT26 reduces global optimality to shortest paths in
the finite type graph (Dijkstra, complete). NT27 proves the fusion bound:
fusing costs no more than composing, and strictly less when the intermediate
type never materializes. NT28 and NT29 formalize refinement and fast-path
specialization: substitution preserves refinement, behavior, and cost
monotonicity, and every finite refinement lattice has a unique most-refined
element per behavior class, so the most-refined dispatcher is sound and
cost-minimal. NT30 closes the complexity class of the dataplane: atoms whose
costs lie in a class closed under the resource operations yield molecules in
that class. NT31 is the descriptive-complexity anchor: finite-state molecules
realize exactly the regular behaviors, and minimal-state molecules are exactly
the minimal DFAs \cite{nerode}\cite{hopcroft}. NT32 is projection reduction:
state that the step function consults only through a projection can be
replaced by the projection's image.

\section{The cost model}\label{sec:enrich-cost}

\begin{definition}[Cost vectors and atom cost assignments]\label{def:cost}
A \emph{cost vector} is an element of $\mathbb{R}_{\ge 0}^{3}$, the three
coordinates being \emph{time}, \emph{memory}, and \emph{cache misses}; write
$c(M) = (t(M), m(M), k(M))$ for the cost of a molecule $M$. The order on cost
vectors is componentwise: $u \le v$ iff $u_i \le v_i$ for $i = 1, 2, 3$, and
$u + v$ is componentwise addition. An \emph{atom cost assignment} for a
signature $\Sigma$ is a function $c \colon \mathrm{At}(\Sigma) \to
\mathbb{R}_{\ge 0}^{3}$ from the atoms of $\Sigma$ to cost vectors.
\end{definition}

\begin{definition}[Additive extension]\label{def:additive-cost}
Given an atom cost assignment, the \emph{additive extension} assigns to every
molecule over $\Sigma$ a cost vector by the clauses
$\mathrm{id}_A \mapsto 0$, $\sigma_{A,B} \mapsto 0$,
$g \circ f \mapsto c(f) + c(g)$, and $f \otimes g \mapsto c(f) + c(g)$.
NT25 shows that this assignment is well-defined (independent of the order in
which the clauses are applied), additive over composition, and subadditive
over tensor.
\end{definition}

\begin{theorem}[Cost enrichment is well-defined]\label{nt:25}
\textbf{NT25 [N].} Fix an atom cost assignment for a signature $\Sigma$
(Definition~\ref{def:cost}). The additive extension
(Definition~\ref{def:additive-cost}) is a well-defined function on the
molecules over $\Sigma$: it is invariant under the structural laws of the free
symmetric monoidal category --- reassociation of composition, insertion and
deletion of identities, and the coherence equations of the braiding --- and it
satisfies, for all molecules $f, g$, the three quantitative laws of a
resource-sensitive monoidal category:
(1) \emph{additive composition cost}: $c(g \circ f) = c(f) + c(g)$;
(2) \emph{triangle inequality}: $c(g \circ f) \le c(f) + c(g)$, with equality,
and more generally $c(f_1 \circ \cdots \circ f_n) \le \sum_{i=1}^{n} c(f_i)$;
(3) \emph{tensor subadditivity}: $c(f \otimes g) \le c(f) + c(g)$, with
equality for the additive extension and strict inequality possible for
implementations that share fixed costs across the two factors (batching,
\autoref{ch:operad}).
Consequently cost is a well-defined resource annotation of Mol: every molecule
has a cost vector computed from its atoms by the inductive rules, independent
of how the molecule is bracketed.
\end{theorem}

\begin{proof}
\emph{Invariance under the structural laws.} Define the cost syntactically: the
cost of a molecule is the sum of the costs of its atoms, identities and
braidings costing zero. This agrees with the clauses of
Definition~\ref{def:additive-cost}, and we show it is invariant under the
symmetric monoidal axioms, checked instance by instance. \emph{Associativity:}
$c((f \circ g) \circ h) = c(f) + c(g) + c(h) = c(f \circ (g \circ h))$, and
similarly for $\otimes$. \emph{Unit laws:} $c(f \circ \mathrm{id}) = c(f) +
0 = c(f)$, and $c(\mathrm{id} \otimes f) = c(f)$. \emph{Braiding:}
$\sigma_{A,B}$ contains no atoms, so $c(\sigma_{A,B}) = 0$, and both sides of
the symmetry law $\sigma_{B,A} \circ \sigma_{A,B} = \mathrm{id}$ and of the
hexagon equations carry cost $0$. \emph{Interchange law:} the two sides of
$(f \otimes g) \circ (h \otimes k) = (f \circ h) \otimes (g \circ k)$ both sum
the costs of $f, g, h, k$. In every case the sum over atoms is unchanged, so
the syntactic cost is invariant under the axioms. Mac Lane's coherence theorem
\cite{maclane} states that these axioms generate exactly the equational theory
of the free symmetric monoidal category, so the cost of a molecule is
independent of the representative chosen, hence well-defined.
\emph{The three laws.} \emph{Additivity} holds because the composite
$g \circ f$ has as its atoms exactly the atoms of $f$ together with those of
$g$: plugging introduces no new nodes. \emph{Triangle inequality.} By
additivity the inequality holds with equality; iterating over a chain of
length $n$ gives the general bound. \emph{Tensor subadditivity.} The tensor
$f \otimes g$ is the disjoint union of the two machines, so the additive
extension gives equality. Any implementation may charge shared setup costs (a
single loop header, one syscall amortization) only once; formally, any cost
assignment satisfying (1)--(3) bounds its tensor component by the additive
one, and the additive extension is the pointwise largest such assignment.
Hence (3) holds in general, with equality in the canonical model.
\end{proof}

\begin{remark}[Cost is defined on molecules, not behavior classes]
Well-definedness is with respect to the structural equations, not behavioral
equivalence: behaviorally equivalent molecules may carry different costs (NT6
only equates behaviors). This is exactly the freedom that refinement
(\autoref{sec:enrich-refine}) exploits: among the realizations of one
behavior, the dataplane may pick the cheapest.
\end{remark}

\section{Global optimality as a shortest-path problem}\label{sec:enrich-shortest}

Cost is additive along composition (NT25), so the cost of a pipeline that is a
sequential composite of atoms is the sum of the atom costs. Optimization is
then graph theory: the atoms are the edges of a finite \emph{type graph}, and
the cheapest pipeline is a shortest path in it.

\begin{theorem}[Global optimality via shortest paths]\label{nt:26}
\textbf{NT26 [N].} Let $\Sigma$ be a finite signature with bounded types and
nonnegative atom costs, additively extended (NT25). The \emph{type graph} $G$
has as vertices the types over $\Sigma$ and, for every atom
$a \colon A \to B$, an edge $A \to B$ of weight $c(a) \ge 0$. For all types
$A, B$:
(1) the minimum cost of a \emph{path-molecule} $A \to B$ --- a molecule
composed of atoms and identities by sequential composition alone --- equals
the shortest-path distance $d(A,B)$ in $G$;
(2) Dijkstra's algorithm computes $d(A,B)$ for every pair of types, terminates
on the finite graph, and is \emph{complete}: whenever $d(A,B) < \infty$ it
returns a shortest path, and that path is a min-cost path-molecule;
(3) if every atom has a single-type source and a single-type target (the
\emph{slice signature} of the dataplane hot path), then every molecule
$A \to B$ is a path-molecule, and (1)--(2) hold for \emph{all} molecules:
globally optimal molecules are exactly the shortest paths of $G$, of cost
$d(A,B)$.
\end{theorem}

\begin{proof}
\emph{(1): the two optimization problems coincide.} A path-molecule $A \to B$
is a finite sequence of atoms $a_1, \dots, a_n$ with $a_i \colon T_{i-1} \to
T_i$, $T_0 = A$, $T_n = B$, interleaved with identities; it denotes the
composite. By additivity (NT25),
$c(a_1 \circ \cdots \circ a_n) = \sum_{i=1}^{n} c(a_i)$, the total weight of
the corresponding path $A \to B$ in $G$. Conversely every path in $G$ is the
atom sequence of a path-molecule. The two sets are in bijection with equal
costs, so the minimum over path-molecules equals the minimum over paths,
$d(A,B)$.
\emph{(2): Dijkstra's algorithm.} Maintain for every vertex a tentative
distance, and repeatedly extract the unvisited vertex of minimal tentative
distance, relaxing its outgoing edges. \emph{Correctness by induction on the
extraction order.} Let $v$ be the vertex extracted at step $k$; all previously
extracted vertices carry their true distances. Any path from the source $s$ to
$v$ leaves the extracted set at some edge $u \to w$ with $u$ extracted and $w$
not yet; its total weight is at least
$\mathrm{dist}(u) + c(u \to w) \ge \mathrm{tentative}(w) \ge
\mathrm{tentative}(v)$, the last by the extraction order, and
$\mathrm{tentative}(v)$ is the weight of an actual path, so
$\mathrm{tentative}(v) = d(s, v)$. Nonnegativity is used exactly at
$\mathrm{dist}(u) + c(u \to w) \ge \mathrm{tentative}(w)$, which holds because
relaxing $u \to w$ set $\mathrm{tentative}(w) \le \mathrm{dist}(u) + c(u \to
w)$ and later relaxations only lower it. The graph is finite (bounded types,
\autoref{ch:free}), so the algorithm terminates after at most $V$ extractions
and $E$ relaxations, $O(E \log V)$ with a priority queue. Recording the
predecessor on each relaxation yields a shortest path to every reachable
target; by (1) that path is a min-cost path-molecule, and running the search
from every source type covers all pairs. This is completeness: every pair with
$d(A,B) < \infty$ is returned with a witness.
\emph{(3): slice signatures.} Assume every atom has a single base type as
source and as target. Let $M \colon A \to B$ be any molecule, realized as a
wiring diagram (\autoref{ch:free}). Every atom has exactly one input port and
one output port, so every port has in-degree and out-degree at most one: the
diagram is a disjoint union of directed paths, each component with exactly one
boundary input and one boundary output. For single types $A$ and $B$ the
boundary is $A \to B$, so there is a single component: a path
$X_0 \to X_1 \to \cdots \to X_n$ with $X_0 = A$ and $X_n = B$, i.e.\ a
path-molecule. Its cost is the weight of that path, at least $d(A,B)$, and by
(2) a shortest path attains $d(A,B)$; hence the global minimum over all
molecules $A \to B$ is exactly $d(A,B)$.
\end{proof}

\begin{remark}[Scope of the statement]
For a general signature, a molecule may route through tensor products: a
parallel composite of paths joined by a braiding. Its cost is then the sum of
its components' costs (NT25), and such parallelism is absorbed into the
signature by batch atoms (NT17, \autoref{ch:operad}, NT47). The theorem's
engineering content is unchanged: with nonnegative additive costs, optimal
composition is a shortest-path problem, and Dijkstra's algorithm is complete
for it. The lookup table of \autoref{ch:free} (NT24) stores exactly the
shortest paths of each behavior class.
\end{remark}

\section{The fusion bound}\label{sec:enrich-fusion}

Cut elimination fuses a pipeline into a single morphism (NT49,
\autoref{ch:linear}); fusion pays for itself. The naive run of $g \circ f$
writes the intermediate value of type $B$ to memory and reads it back; the
fused run keeps that value in registers.

\begin{definition}[Measured cost of materialization]\label{def:measured}
Let $m(B) \ge 0$ be the memory charge of a materialized intermediate value of
type $B$, with $m(B) > 0$ whenever $B$ is nonempty. The \emph{measured cost}
of a naive run of the composite $g \circ f$ is
$c^{\mathrm{meas}}(g \circ f) = c(f) + c(g) + m(B)$: the atoms do their nominal
work and the intermediate is stored and reloaded. The \emph{canonical fused
molecule} for the pair $(f, g)$ is the same atoms rewired so that the
intermediate is consumed in registers; it realizes $g \circ f$ and its
measured cost is $c^{\mathrm{meas}}(h) = c(f) + c(g)$: no materialization
charge. Any \emph{fused realization} of the behavior is a molecule $h$ with
$h \simeq g \circ f$ and measured cost at most the canonical one.
\end{definition}

\begin{theorem}[Fusion bound]\label{nt:27}
\textbf{NT27 [N].} Let $f \colon A \to B$ and $g \colon B \to C$ be molecules.
(1) The naive composite realizes $g \circ f$ at measured cost
$c(f) + c(g) + m(B)$; the canonical fused molecule realizes the same behavior
at cost $c(f) + c(g)$. Hence the optimal fused realization $h^{*}$ satisfies
the \emph{fusion bound} $c(h^{*}) \le c(f) + c(g)$.
(2) The bound is strict exactly when the intermediate type never materializes:
if $B$ is nonempty, then $m(B) > 0$ and
$c^{\mathrm{meas}}(h^{*}) \le c(f) + c(g) < c(f) + c(g) + m(B) =
c^{\mathrm{meas}}(g \circ f)$ --- fusing strictly decreases the measured cost,
and in particular the memory and cache-miss components improve by the
eliminated store and load. If $B$ is empty, $m(B) = 0$ and the two coincide.
\end{theorem}

\begin{proof}
\emph{(1): the canonical fused molecule exists and realizes $g \circ f$.} The
composite $g \circ f$ has state space $S_f \times S_g$: first $f$ transforms
the input and updates its state, then $g$ consumes the intermediate and
updates its state. The canonical fused molecule is the same state space and
the same two steps, with the intermediate held in registers instead of memory;
its atom multiset is exactly that of $g \circ f$ (plugging introduces no
atoms, NT25), so its nominal cost is $c(f) + c(g)$ and, since it never
materializes $B$, its measured cost is $c(f) + c(g)$
(Definition~\ref{def:measured}). It realizes $g \circ f$ by construction,
while the naive run additionally stores and reloads the intermediate, costing
$c(f) + c(g) + m(B)$. The optimal fused realization is no more expensive than
the canonical one, giving $c(h^{*}) \le c(f) + c(g)$.
\emph{(2): strictness.} If $B$ is nonempty then $m(B) > 0$, so
$c(f) + c(g) < c(f) + c(g) + m(B)$: the left side is the canonical fused cost,
the right side the naive composite cost, and fusing strictly decreases the
measured cost. In the memory and cache-miss components the difference is
exactly the charge of the eliminated store and load, so those components
improve strictly. If $B$ is the empty type, $m(B) = 0$, both runs cost
$c(f) + c(g)$, and the bound is tight.
\end{proof}

\begin{remark}[Fusion, normalization, and the lookup table]
Fusion is cost-decreasing, and the normalization of \autoref{ch:rewrite} is
cost-nonincreasing (NT18). Together with the fusion bound they imply that the
normal form of a composite --- the fused pipeline --- is a refinement of every
enumerated representative of its behavior class; this is the mechanism by
which the lookup table of \autoref{ch:free} stores cost-optimal molecules
(NT24, NT26, NT29).
\end{remark}

\section{Refinement and fast-path specialization}\label{sec:enrich-refine}

Inside a behavior class, molecules are ordered by cost. A \emph{refinement}
replaces a molecule by a cheaper one with the same behavior --- the formal
counterpart of fast-path specialization: the general path is correct, and the
specialized path is also correct and faster.

\begin{definition}[Refinement]\label{def:refine}
For molecules $f, f' \colon A \to B$ with $f \simeq f'$, write
$f \sqsubseteq f'$ and call $f'$ a \emph{refinement} of $f$ when
$c(f') \le c(f)$ componentwise. Write $f \equiv f'$ when both
$f \sqsubseteq f'$ and $f' \sqsubseteq f$ hold (equal behavior and equal
cost). The relation $\sqsubseteq$ is a preorder on each behavior class, and a
partial order modulo $\equiv$.
\end{definition}

\begin{theorem}[Refinement substitution]\label{nt:28}
\textbf{NT28 [N].} Refinement is preserved by substitution into compositions.
Let $f, f' \colon A \to B$ and $g, g' \colon B \to C$ be molecules with
$f \sqsubseteq f'$ and $g \sqsubseteq g'$. Then $g \circ f \sqsubseteq
g' \circ f'$: the composite $g' \circ f'$ is a refinement of $g \circ f$.
Behavior is preserved --- $g' \circ f' \simeq g \circ f$ --- and cost is
monotone --- $c(g' \circ f') \le c(g \circ f)$. The same holds for the tensor:
$f \otimes g \sqsubseteq f' \otimes g'$ whenever the tensor cost is additive on
the pairs involved (in particular under the additive extension, NT25).
\end{theorem}

\begin{proof}
\emph{Behavior preserved.} From $f \sqsubseteq f'$ we have $f \simeq f'$, and
from $g \sqsubseteq g'$ we have $g \simeq g'$. Behavioral equivalence is a
congruence for $\circ$ and $\otimes$ (NT6), so $g \circ f \simeq g' \circ f'$
and $f \otimes g \simeq f' \otimes g'$.
\emph{Cost monotone, sequential.} By additivity (NT25),
$c(g \circ f) = c(f) + c(g)$ and $c(g' \circ f') = c(f') + c(g')$. Since
$c(f') \le c(f)$ and $c(g') \le c(g)$ componentwise, adding componentwise
gives $c(f') + c(g') \le c(f) + c(g)$, i.e.\ $c(g' \circ f') \le
c(g \circ f)$. Together with the congruence, $g \circ f \sqsubseteq
g' \circ f'$.
\emph{Cost monotone, tensor.} Under the additive extension,
$c(f \otimes g) = c(f) + c(g)$ and $c(f' \otimes g') = c(f') + c(g')$ (NT25),
and the same inequality gives $c(f' \otimes g') \le c(f \otimes g)$; with the
tensor congruence, $f \otimes g \sqsubseteq f' \otimes g'$. For an arbitrary
cost assignment the same conclusion holds whenever the tensor cost equals the
sum on both pairs, which is the stated hypothesis.
\end{proof}

\begin{remark}[Fast-path specialization]
The substitution property is the soundness certificate of specialization: any
component of a pipeline may be replaced, at any depth, by a refinement of
itself, and the whole pipeline remains a refinement of the original --- same
behavior, no greater cost. This is how the dataplane's fast paths are
introduced without re-verifying the full pipeline: each specialized atom is
checked once against the general one (behavioral equivalence, NT6 and
\autoref{ch:testing}), and NT28 lifts the guarantee to every composite.
\end{remark}

\begin{theorem}[Maximal refinement]\label{nt:29}
\textbf{NT29 [N].} Let $\mathcal{C}$ be a finite set of molecules
$A \to B$ with a common behavior (a \emph{behavior class}), ordered by
$\sqsubseteq$ modulo $\equiv$ (Definition~\ref{def:refine}).
(1) \emph{Existence.} $\mathcal{C}$ has a maximal element --- a molecule
refined by no strictly better molecule of $\mathcal{C}$ --- and every element
of $\mathcal{C}$ is refined by some maximal element.
(2) \emph{Uniqueness.} If $\mathcal{C}$ is \emph{cost-closed} --- it contains
a molecule whose cost is the componentwise minimum of the costs occurring in
$\mathcal{C}$ --- then the maximal element is unique up to $\equiv$: the
\emph{most-refined} representative, of cost $\min_{f \in \mathcal{C}} c(f)$.
(3) \emph{The most-refined dispatcher}, which selects for every behavior class
its most-refined representative, is sound --- each selected molecule has the
class behavior, because it lies in the class --- and cost-minimal --- no other
selection realizes the class at strictly smaller cost in any component.
\end{theorem}

\begin{proof}
\emph{(1): maximal elements exist.} If $f \in \mathcal{C}$ is not maximal, some
$f_1 \in \mathcal{C}$ strictly refines it: $c(f_1) \le c(f)$ componentwise and
$c(f_1) \neq c(f)$. A strictly refining chain
$f, f_1, f_2, \dots$ cannot be infinite: each step lowers at least one
component, and the componentwise-decreasing sequences in the finite set of
cost vectors of $\mathcal{C}$ are finite. Hence a maximal element is reached,
and every element is refined by one (iterate the same argument from it).
\emph{(2): uniqueness under cost-closure.} Let
$v^{*} = \min_{f \in \mathcal{C}} c(f)$ componentwise; by cost-closure some
$f^{*} \in \mathcal{C}$ has $c(f^{*}) = v^{*}$. Let $g$ be any maximal element
of $\mathcal{C}$. Since $v^{*} \le c(g)$ componentwise, if $c(g) \neq v^{*}$
then $f^{*}$ strictly refines $g$, contradicting maximality; hence
$c(g) = v^{*}$. All maximal elements therefore carry the cost $v^{*}$ and the
class behavior, so they are pairwise $\equiv$; and $f^{*}$ itself is maximal,
so maximal elements exist and are unique up to $\equiv$.
\emph{(3): soundness and cost-minimality of the dispatcher.} The most-refined
representative of a class is an element of the class, so it realizes the class
behavior: the dispatcher is sound. Its cost is $v^{*} \le c(f)$ for every
$f$ in the class, so no other selection realizes the class at strictly
smaller cost in any component: the dispatcher is cost-minimal.
\end{proof}

\begin{remark}[Tradeoffs and decision matrices]
Without cost-closure there can be several incomparable maximal refinements ---
one fast but memory-heavy, one slow but memory-light --- and this is exactly
the tradeoff that the decision matrices of \autoref{ch:situations} resolve
with a weight vector over the cost components. For the dataplane,
cost-closure is supplied by the lookup table of \autoref{ch:free}:
normalization is cost-nonincreasing (NT18), so the normal form of a class is a
refinement of every enumerated representative (NT22--NT24), and fusion (NT27)
only improves it further.
\end{remark}

\section{Complexity closure}\label{sec:enrich-closure}

The complexity of the dataplane is decided by its atoms: whatever resource
class the atoms belong to, and which the two compositions preserve, belongs to
every molecule built from them.

\begin{theorem}[Complexity closure]\label{nt:30}
\textbf{NT30 [N].} (1) \emph{Structural closure.} Let $P$ be a property of
molecules preserved by the composition operations --- if $f$ and $g$ have $P$
then $g \circ f$ and $f \otimes g$ have $P$, and all identities and braidings
have $P$. If every atom of a signature $\Sigma$ has $P$, then every molecule
over $\Sigma$ has $P$.
(2) \emph{Resource closure.} Let $F$ be a class of nonnegative cost bounds
closed under addition (covering the time and cache-miss components), and let
$G$ be a class of state-size bounds closed under multiplication and containing
$1$ (covering the memory component, because state spaces multiply under both
compositions). If every atom of $\Sigma$ has time and cache-miss cost in $F$
and state size in $G$, then every molecule over $\Sigma$ has time and
cache-miss cost in $F$ and state size in $G$.
\end{theorem}

\begin{proof}
\emph{(1): induction on the construction of molecules.} Every molecule over
$\Sigma$ is built from atoms, identities, and braidings by finitely many
applications of $\circ$ and $\otimes$: this is the definition of a molecule
(NT3, \autoref{ch:mol}), and the free-category construction of
\autoref{ch:free} shows that the construction is well-founded. Induct on the
number of operations. Base cases: an atom has $P$ by hypothesis, and the
identities and braidings have $P$ by hypothesis. Induction steps: if $f$ and
$g$ have $P$, then $g \circ f$ and $f \otimes g$ have $P$ because $P$ is
preserved by both operations. Hence every molecule over $\Sigma$ has $P$.
\emph{(2): resource bounds by induction, using the composition laws of NT25.}
For the time and cache-miss components, additivity gives
$t(g \circ f) = t(f) + t(g)$ and $k(g \circ f) = k(f) + k(g)$, and tensor
subadditivity gives $t(f \otimes g) \le t(f) + t(g)$ and
$k(f \otimes g) \le k(f) + k(g)$; since $F$ is closed under addition, the
components stay in $F$ at every induction step. For the memory component, the
state space of $g \circ f$ is $S_f \times S_g$, and the state space of
$f \otimes g$ is likewise $S_f \times S_g$, so
$|S_{g \circ f}| = |S_f| \cdot |S_g| = |S_{f \otimes g}|$; the identity has
$|S_{\mathrm{id}}| = 1 \in G$, and by induction over the construction the
state sizes remain in $G$, which is closed under multiplication. This proves
the resource closure.
\end{proof}

\begin{remark}[The bounded-resource consequence]
For $F$ the bounded functions (atoms with bounded time) and $G = \{1\}$ up to
a constant (atoms with bounded state), closure gives: every molecule has
bounded time per input and bounded state --- subject to the same constants
raised by the multiplicative composition of state bounds. The zero-allocation
discipline of \autoref{ch:linear} (NT50) is the engineering form of this
statement: the complexity class of the dataplane is the complexity of its
atoms, never more.
\end{remark}

\section{Descriptive complexity}\label{sec:enrich-descriptive}

Finite-state molecules sit exactly at the regular level of the Chomsky
hierarchy. Fix an input alphabet $A$ and an output alphabet $B$ for a
finite-state molecule $M = (S, \mathit{step})$, an initial state $s_0 \in S$,
and a set $F \subseteq S$ of accepting states. The \emph{language recognized}
by $M$ is $L(M) = \{w \in A^{*} : \text{the run of } M \text{ on } w \text{ ends
in } F\}$; the \emph{transduction computed} by $M$ from $s_0$ is the
input--output map of \autoref{ch:mol}.

\begin{theorem}[Descriptive complexity of Mol]\label{nt:31}
\textbf{NT31 [C].} (1) \emph{Regularity.} Every finite-state molecule
recognizes a regular language, and its transduction is a deterministic Mealy
transduction \cite{mealy}: finite state realizes no behavior beyond the
regular ones.
(2) \emph{Completeness.} Every regular language is recognized by some
finite-state molecule, namely the minimal DFA of the language, whose states
are the Myhill--Nerode equivalence classes \cite{nerode}.
(3) \emph{Minimality.} Among the finite-state molecules recognizing a regular
language $L$, the minimum number of states is the number $\kappa(L)$ of Nerode
classes of $L$; every molecule attaining it --- a \emph{minimal-state
molecule} --- is a minimal DFA, unique up to state-space isomorphism, and
Hopcroft's algorithm computes it in time $O(n \log n)$ from any $n$-state
realization \cite{hopcroft}.
\end{theorem}

\begin{proof}
\emph{(1): finite state stays inside the regular languages.} Write
$\delta(s, a)$ for the state component of $\mathit{step}(s, a)$. The molecule
is a finite deterministic automaton $(S, A, \delta, s_0, F)$, and every such
automaton recognizes a regular language. Concretely, by state elimination:
let $L_s \subseteq A^{*}$ be the language recognized from state $s$. The
languages $L_s$ satisfy the finite system of equations
\[
L_s = \bigcup_{a \in A} a \cdot L_{\delta(s,a)},
\]
with the term $\varepsilon$ added when $s \in F$. Solving the system by
substituting equations into each other --- the derivative method of Brzozowski
\cite{brzozowski} --- expresses every $L_s$, in particular $L_{s_0} = L(M)$,
as a regular expression over $A$; hence $L(M)$ is regular. The transduction
reading is literal: a finite-state molecule with deterministic step is a
deterministic Mealy machine (NT8, \autoref{ch:mol}), so its transduction is a
deterministic Mealy transduction \cite{mealy}.
\emph{(2): every regular language is recognized by a finite-state molecule.}
Let $L \subseteq A^{*}$ be regular. The Nerode equivalence $u \equiv_L v$ iff
$\forall w \in A^{*} \colon (uw \in L \Leftrightarrow vw \in L)$ has finitely
many classes exactly when $L$ is regular \cite{nerode}. The quotient
$S_L = A^{*} / \equiv_L$ with transition $\delta([u], a) = [ua]$ and accepting
set $\{[u] : u \in L\}$ is a finite deterministic automaton recognizing $L$,
well-defined because the class definition makes the transition and the
accepting predicate independent of the representative. It is a finite-state
molecule with state space $S_L$ and $\mathit{step}([u], a) = (o, [ua])$ for
any fixed $o \in B$, the accepting set carried separately, and $L(M) = L$.
\emph{(3): minimal states are Nerode classes.} \emph{Lower bound.} Any
automaton recognizing $L$ has at least $\kappa(L)$ states: words $u, v$ with
$u \not\equiv_L v$ reach states that must be distinct, since one state cannot
both accept and reject the distinguishing continuation $w$.
\emph{Upper bound and uniqueness.} The automaton of (2) has exactly
$\kappa(L)$ states, so the minimum is $\kappa(L)$. Let $N$ be any
$\kappa(L)$-state automaton recognizing $L$ and let $\mathrm{st}_N(u)$ be the
state it reaches on $u$. The map $\phi \colon \mathrm{st}_N(u) \mapsto [u]$ is
well-defined: if $u$ and $v$ reached the same state, that state cannot
separate $uw$ from $vw$, so $u \equiv_L v$. It is injective ($u \not\equiv_L
v$ forces $\mathrm{st}_N(u) \neq \mathrm{st}_N(v)$, by the lower-bound
argument) and surjective (both sets have $\kappa(L)$ elements). It preserves
the initial state ($\phi(\mathrm{st}_N(\varepsilon)) = [\varepsilon]$), the
transitions ($\phi(\mathrm{st}_N(ua)) = [ua] = \delta_L([u], a)$), and the
accepting sets ($\mathrm{st}_N(u) \in F_N \Leftrightarrow u \in L
\Leftrightarrow [u] \in F_L$). Hence $\phi$ is an isomorphism of machines, and
minimal-state molecules are unique up to state-space isomorphism; by the
state-space bijection quotient of \autoref{ch:mol} (NT2) the minimal molecule
is unique. \emph{Complexity.} Hopcroft's algorithm computes the coarsest
stable partition of the states of a given $n$-state automaton --- exactly the
partition into Nerode classes --- in time $O(n \log n)$ \cite{hopcroft};
applied to any $n$-state realization of $L$, it returns the minimal molecule.
\end{proof}

\begin{remark}[The transduction reading]
For transductions the same statement holds with the Nerode relation replaced
by the behavioral equivalence $\simeq$ of \autoref{ch:mol}: the minimal-state
molecule realizing a transduction is the quotient of any realization by
$\simeq$, unique up to state-space isomorphism (NT38, \autoref{ch:coalgebra}).
The theorem shows that for the language projection this quotient is exactly
the Myhill--Nerode quotient, so the two minimality theories coincide.
\end{remark}

\section{Projection reduction}\label{sec:enrich-projection}

Real molecules carry state that matters only through part of itself. When the
step function reads the state only through a projection, the state space can
be replaced by the projection's image without changing behavior and without
losing minimality.

\begin{theorem}[Projection reduction]\label{nt:32}
\textbf{NT32 [N].} Let $M = (S, \mathit{step})$ be a finite-state molecule and
let $\pi \colon S \to P$ be a surjection such that the step function factors
through $\pi$: there is a function
$\mathit{step}_{P} \colon P \times A \to B \times P$ with
$\mathit{step}(s, a) = \mathit{step}_{P}(\pi(s), a)$ for all $s \in S$ and
$a \in A$ --- the step depends on the state only through $\pi$. Then:
(1) $(P, \mathit{step}_{P})$ is a well-defined finite-state molecule, the
\emph{projection reduction} of $M$;
(2) for every initial state $s_0$, the reduced molecule started at
$\pi(s_0)$ is behaviorally equivalent to $M$ started at $s_0$;
(3) if $\pi$ is \emph{discriminating} --- no two distinct elements of $P$ are
behaviorally equivalent under $\mathit{step}_{P}$ --- then the minimal state
space of the behavior satisfies $S_{\min} \cong \pi(S) = P$: the projection
reduction is already the minimal realization.
\end{theorem}

\begin{proof}
\emph{(1): the reduction is a molecule.} The factorization says exactly that
$\mathit{step}$ is constant on the fibers of $\pi$: if $\pi(s) = \pi(s')$,
then $\mathit{step}(s, a) = \mathit{step}_{P}(\pi(s), a) =
\mathit{step}_{P}(\pi(s'), a) = \mathit{step}(s', a)$. Hence
$\mathit{step}_{P}$ is well-defined on $P$, and $(P, \mathit{step}_{P})$ is a
finite-state molecule in the sense of \autoref{ch:mol}.
\emph{(2): behavioral equivalence.} Let $(s_i)_{i \ge 0}$ be the run of $M$
on the word $a_1 \cdots a_n$ from $s_0$, and $(p_i)_{i \ge 0}$ the run of the
reduction from $\pi(s_0)$. We show $p_i = \pi(s_i)$ by induction: for $i = 0$
this is the initialization, and if $p_i = \pi(s_i)$, the next reduced state is
the state component of
$\mathit{step}_{P}(p_i, a_{i+1}) = \mathit{step}_{P}(\pi(s_i), a_{i+1}) =
\mathit{step}(s_i, a_{i+1})$, which is $\pi(s_{i+1})$; the outputs coincide at
every step by the same equality. Hence the two machines emit the same output
word on every input word: they are behaviorally equivalent (NT6).
\emph{(3): minimality when $\pi$ is discriminating.} Two states in the same
fiber of $\pi$ are behaviorally equivalent: by (2) applied to the sub-machines
started at $s$ and $s'$ with $\pi(s) = \pi(s')$, their runs agree forever.
Hence every behavioral equivalence class is a union of fibers, and the minimal
state space $S_{\min}$ --- the quotient of $S$ by behavioral equivalence
(NT38, \autoref{ch:coalgebra}) --- is a quotient of the fiber quotient
$P = S / \ker \pi$: a canonical surjection $P \twoheadrightarrow S_{\min}$
gives $|S_{\min}| \le |P|$. Since $\pi$ is discriminating, the elements of
$P$ realize pairwise distinct behaviors, so any realization needs at least
$|P|$ states; the reduction provides exactly $|P|$, hence $|S_{\min}| =
|P|$. The canonical surjection is then a bijection, and it is an isomorphism
of machines: both sides send $\pi(s)$ to its behavioral equivalence class, and
the step is preserved by construction. Thus $S_{\min} \cong P = \pi(S)$; by
the state-space bijection quotient of Mol (NT2) the isomorphism is an equality
of Mol-morphisms.
\end{proof}

\begin{remark}[State hygiene in the dataplane]
The step functions of real dataplane atoms consult a handful of fields of a
large state record. Projection reduction is the formal license to drop the
rest: the molecule that carries only the consulted projection is
behaviorally identical and, when the projection is discriminating, already
minimal --- the layout discipline of policy [CACHE] of the companion standard.
\end{remark}

\section{Summary}\label{sec:enrich-summary}

\begin{table}[htbp]
\centering
\begin{tabular}{@{}lll@{}}
\toprule
\textbf{Theorem} & \textbf{Tag} & \textbf{Content} \\
\midrule
NT25 & [N] & cost enrichment well-defined: additivity, triangle inequality, tensor subadditivity \\
NT26 & [N] & optimal molecules are shortest paths in the type graph; Dijkstra complete \\
NT27 & [N] & fusion bound $c(h^{*}) \le c(f) + c(g)$; strict when the intermediate never materializes \\
NT28 & [N] & refinement substitution; behavior preserved, cost monotone \\
NT29 & [N] & unique most-refined element per class; dispatcher sound and cost-minimal \\
NT30 & [N] & complexity closure under $\circ$ and $\otimes$ \\
NT31 & [C] & finite-state molecules = regular behaviors; minimal states = minimal DFA \\
NT32 & [N] & projection reduction $S_{\min} \cong \pi(S)$ \\
\bottomrule
\end{tabular}
\caption{The theorems of this chapter.}
\label{tab:enrich}
\end{table}

This chapter gives Mol its quantitative layer. Costs form a well-defined
enrichment of the category (NT25): composition adds, tensor is subadditive,
and the triangle inequality holds. With nonnegative additive costs, global
optimality is a shortest-path problem in the finite type graph, and Dijkstra's
algorithm is complete for it (NT26). Fusion always pays: the fused composite
costs no more than the sum of its parts, and strictly less when the
intermediate type never materializes (NT27). Refinement gives fast-path
specialization a soundness certificate: substitution preserves behavior and
monotonically improves cost (NT28), and the most-refined dispatcher is sound
and cost-minimal (NT29). Complexity closes over composition: the resource
class of the dataplane is the resource class of its atoms (NT30). The
expressive ceiling is exactly the regular level of the hierarchy (NT31), and
projection reduction prunes state that is only consulted through a projection
(NT32). Together with the normal forms of \autoref{ch:free} and the decision
matrices of \autoref{ch:situations}, these results turn ``choose a molecule
for this behavior within this budget'' into a finite, decidable, and
computable problem (NT51).
